\begin{prob}
    In each of the parts below, compute the \emph{expected running time} of the
    given code. Also state the \emph{worst-case} running time. Use asymptotic
    notation.
    
    For full credit, for each part:
    (1) clearly state the cases you consider,
    (2) give the probability of each case,
    (3) give the time taken in each case, and
    (4) compute the expected time as a probability-weighted sum.
    
    \textbf{Reminder.} In Python, \python{np.random.randint(n)} picks a random
    integer uniformly from $\{0,1,\ldots,n-1\}$.
    
    \begin{subprobset}
    
    \begin{subprob}
        \begin{minted}[autogobble]{python}
            def alpha(n):
                # pick a random integer from 0 to n-1
                r = np.random.randint(n)
                if r < n/5:
                    for i in range(n**3):
                        print(i)
                else:
                    for i in range(n):
                        print(i)
        \end{minted}
    \end{subprob}
    
    \begin{subprob}
        \begin{minted}[autogobble]{python}
            def beta(n):
                # pick a random integer from 0 to n-1
                r = np.random.randint(n)
                if r < 20:
                    for i in range(n**2):
                        print(i)
                else:
                    for i in range(n):
                        print(i)
        \end{minted}
    \end{subprob}
    
    \begin{subprob}
        \begin{minted}[autogobble]{python}
            def gamma(n):
                # pick a random integer from 0 to n-1
                m = np.random.randint(n)
                total = 0
                for i in range(m):
                    for j in range(i):
                        total += 1
        \end{minted}
    \end{subprob}
    
    \begin{subprob}
        \begin{minted}[autogobble]{python}
            def random_search(arr, target):
                n = len(arr)
                while True:
                    idx = np.random.randint(n)
                    if arr[idx] == target:
                        return idx
        \end{minted}
        Assume all elements in \python{arr} are distinct and \python{target} is in the array.
    \end{subprob}
    
    \end{subprobset}
    
    \begin{soln}
    \begin{itemize}
        \item \textbf{(a)}  
        There are two cases.  
        If $r < n/5$ (probability $1/5$), the loop runs $n^3$ times, so the time is
        $\Theta(n^3)$.  
        Otherwise (probability $4/5$), the loop runs $n$ times, so the time is
        $\Theta(n)$.  
    
        The expected time is
        \[
        \mathbb{E}[T] = \frac{1}{5}\Theta(n^3) + \frac{4}{5}\Theta(n) = \Theta(n^3).
        \]
        The worst-case time complexity is $\Theta(n^3)$.
    
        \item \textbf{(b)}  
        If $r < 20$ (probability $20/n$), the running time is $\Theta(n^2)$.  
        Otherwise (probability $1 - 20/n$), the running time is $\Theta(n)$.  
    
        The expected time is
        \[
        \mathbb{E}[T]
        = \frac{20}{n}\Theta(n^2) + \left(1-\frac{20}{n}\right)\Theta(n)
        = \Theta(n).
        \]
        The worst-case time complexity is $\Theta(n^2)$.
    
        \item \textbf{(c)}  
        For a fixed value $m = k$, the inner loop executes
        $1 + 2 + \cdots + (k-1) = \Theta(k^2)$ times.
        Since $m$ is chosen uniformly from $\{0,\ldots,n-1\}$, each case has
        probability $1/n$.
    
        The expected time is
        \[
        \mathbb{E}[T]
        = \sum_{k=0}^{n-1} \frac{1}{n}\Theta(k^2)
        = \frac{1}{n}\Theta(n^3)
        = \Theta(n^2).
        \]
        The worst-case time complexity is $\Theta(n^2)$.
    
        \item \textbf{(d)}  
        On each iteration, the probability of success is $1/n$ and the probability
        of failure is $1 - 1/n$.
        Let Case $\alpha$ be the event that the algorithm succeeds on the
        $\alpha$-th iteration.
        Then
        \[
        \Pr[\text{Case }\alpha] = (1 - 1/n)^{\alpha-1}\cdot \frac{1}{n},
        \]
        and the time taken in Case $\alpha$ is $\Theta(\alpha)$.
    
        The expected running time is therefore $\Theta(n)$.
        The worst-case running time is unbounded.
    \end{itemize}
    
    \textbf{Key takeaway:} Expected running time is computed by identifying all
    possible cases, determining their probabilities, and summing the running time
    of each case weighted by its probability.
    \end{soln}
    
    \end{prob}
